\section{Triangles in Random Graphs}

\Cref{Ch1:Thm:Erdos_Girth} put the random graph $G_{n, p}$ of \Cref{Ch1:Def:Random_Graph} to work building a graph with no short cycles, where all that mattered was that short cycles were scarce. The question here is the opposite one, and a sharper one: for which $p$ does $G_{n, p}$ contain a triangle at all?

Let $X$ be the number of triangles in $G_{n, p}$. We get
\begin{align*}
    \E(X) &= \binom{n}{3} p^3
\end{align*}

% [CORRECTED] was: "so $G_{n, p}$ is likely to have a triangle" --- a diverging expectation does not on its own
% make the event likely ($X = n^2$ with probability $\frac{1}{n}$ has $\E(X) \to \infty$ and $\Pr\of{X \geq 1} \to 0$),
% and for subgraph counts the threshold is set by the densest subgraph rather than by where the expectation
% crosses $1$. The conclusion is true for triangles, but it needs the variance, which is what the rest of the
% section sets up.
If $p = \frac{\omega(1)}{n}$ then $\E(X) \to \infty$, which suggests that $G_{n, p}$ has a triangle.

If $p = o\of{\frac{1}{n}}$, then $\E(X) \to 0$, so the probability that $G_{n, p}$ has a triangle goes to $0$; this direction is \Cref{Ch3:Thm:Markov} at $a = 1$.

None of this actually tells us whether $G_{n, p}$ actually has a triangle or not. We need something more sophisticated than the expectation, so we turn to the variance.

\Cref{Ch3:Cor:Second_Moment} is what the first regime wants: show $\frac{\Var{X}}{\E(X)^2} \to 0$ when $p = \frac{\omega(1)}{n}$, and $G_{n, p}$ has a triangle with probability tending to $1$. That leaves $\Var{X}$ for the triangle count to be computed. \sorry
